% Preview probe: glyphs whose program, encoding or transform is not plain
% Computer Modern: virtual fonts over 8r-encoded base fonts (Times, T1
% ligatures and kerns), SlantFont and ExtendFont from the font map, font
% expansion and letter spacing (microtype), and Latin Modern in T1. One
% kind per page.
\documentclass{article}
\usepackage[T1]{fontenc}
\usepackage[expansion=true,protrusion=true,letterspace=200]{microtype}
\pagestyle{empty}
\newcommand\words{Affine fluffy office -- --- ``quoted'' AVAW To. Ty. f\/i fi ffi ffl}
\begin{document}
\fontfamily{ptm}\selectfont
\noindent Times through its T1 virtual font: \words\par
\noindent\textit{Italic: \words}\par
\noindent\textsl{Slanted (SlantFont .167 in the map): \words}\par
\noindent\textsc{Small caps (a VF over the 8r font)}
\newpage

\fontfamily{phv}\fontseries{mc}\selectfont
\noindent Helvetica narrow (ExtendFont .82 in the map): \words
\newpage

\fontfamily{ptm}\selectfont
\noindent Expansion: a paragraph set by microtype with font expansion, so
that each line's glyphs are drawn from a font expanded or contracted by a
few thousandths, with protrusion of the punctuation into the margins as
well. It runs for a few lines so that the stretch differs from line to
line, which only shows if every glyph is at its place and of its width.
More text, more lines, more expansion, until the paragraph has enough of
them to show it, and a short last line.
\newpage

\noindent Letter spacing: \textls{Spaced Out Letters} and \textls[-50]{tight}.
\newpage

\fontfamily{lmr}\selectfont
\noindent Latin Modern in T1: \words. Accents: \'e \`a \^o \"u \c{c} \v{s} \aa{} \ss.
\end{document}
